\documentclass[11pt,a4paper,twoside,openright]{report}
%\usepackage{hyperref}
\usepackage[pdftex]{graphicx} % Biblioteca para uso de figuras
\usepackage{color}

\usepackage[brazil]{babel} % Biblioteca para uso da l�ngua portuguesa
\usepackage[T1]{fontenc} % Biblioteca para uso da acentua��o de entrada
\usepackage[latin1]{inputenc} % Biblioteca para uso da acentua��o de sa�da

\usepackage{amsthm,amsfonts,amsmath,amssymb}  % Biblioteca para uso de comandos matem�ticos
\usepackage{pslatex}
\usepackage{pstricks,pst-node,color,pst-gantt,pst-coil}
\usepackage{scalefnt}
\usepackage{float} %Permite colocar "\begin{figure}[H]" e colocar imagem exatamente onde desejar
\usepackage[hyphens]{url} % Para Aceitar URL nas referencias
\usepackage{pdfpages}


\usepackage{eurosym} %Pacote para possibilitar o uso do s�mbolo de euro "\euro"

\usepackage{listings} % para importa��o de c�digos fonte 

\usepackage[pdftex]{hyperref}
\hypersetup{bookmarks = true, % mostra a barra de bookmarks
	pdftitle = {O titulo do seu Trabalho}, % titulo
	pdfauthor = {Evandro L. L. Rodrigues}, % autor
	pdfsubject = {TCC}, % subject of the document
	pdfkeywords = {Sistemas Embarcados, ARM, Linux Embarcado}, % keywords
	colorlinks = true, 
	linkcolor = black,  
	citecolor = black,
	urlcolor = blue
}


\usepackage{comment}
\usepackage[margin=2.7cm]{geometry}
\renewcommand{\baselinestretch}{1.5}

\setlength{\parskip}{0em}

% Pacote para configurar cabe�alho e rodape
\usepackage{fancyhdr}
\pagestyle{empty}
\fancyhf{} % clear all header and footer fields

\fancypagestyle{plain}{\pagestyle{fancy}}
\renewcommand{\headrulewidth}{0pt}
\renewcommand{\footrulewidth}{0pt}


%Pacote para organizar apendices 
\usepackage[titletoc,toc]{appendix}
\renewcommand{\appendixtocname}{Ap�ndices}


\newenvironment{poliabstract}[1]
  {\renewcommand{\abstractname}{#1}\begin{abstract}}
  {\end{abstract}}

%%%%%%%%%%%%%%%%%%%%%%%%%%%%%%%%%%%%%% CONFIGURA��ES DE P�GINA %%%%%%%%%%%%%%%%%%%%%%%%%%%%%%%%%%%%%%
%\topmargin -2.1cm
%\oddsidemargin 0.5cm 
%\evensidemargin 0.5cm 
%\textwidth 15cm
%\textheight 25.1cm
	
%%%%%%%%%%%%%%%%%%%%%%%%%%%%%%%%%%%%%%%% IN�CIO DO DOCUMENTO %%%%%%%%%%%%%%%%%%%%%%%%%%%%%%%%%%%%%%%%
\begin{document}
	
%%%%%%%%%%%%%%%%%%%%%%%%%%%%%%%%%%%%%%  INCLUDES %%%%%%%%%%%%%%%%%%%%%%%%%%%%%%%%%%%%%%

%%%%%%%%%%%%%%%%%%%%%%%%%%%%%%%%%%%%%% ELEMENTOS PR�-TEXTUAIS %%%%%%%%%%%%%%%%%%%%%%%%%%%%%%%%%%%%%%
%%%%%%%%%%%%%%%%%%%%%%%%%%%%%%%%%%%%%%% CONFIGURA��ES DE CAPA %%%%%%%%%%%%%%%%%%%%%%%%%%%%%%%%%%%%%%%
\begin{titlepage}
	
	% CAPA PRINCIPAL
	\begin{center}
		\Huge{UNIVERSIDADE DE S�O PAULO}\\
		\vspace{0.02\textheight}
		\huge{ESCOLA DE ENGENHARIA DE S�O CARLOS}\\
		\vspace{0.01\textheight}
		\huge{DEPTO. DE ENGENHARIA EL�TRICA E DE COMPUTA��O}\\
		\vspace{0.2\textheight}
		\huge{\textbf{SEL0630 - Aplica��o de Microprocessadores II}}\\
		\vspace{0.05\textheight}
		\huge{\textbf{Proposta de projeto - Monitoramento remoto inteligente}}
		\vspace{0.05\textheight}
	\end{center}
		{
			\begin{flushleft}
			\Large{ \textbf{Autor}: \hspace{1cm} Gabriel Arantes de Lima e Silva, n$^o$. USP 9266372 }\\
			\Large{ \textbf{Autor}: \hspace{1cm} Gabriel Martins da Silva, n$^o$. USP 9266747 }\\
			\Large{ \textbf{Orientador}: \hspace{0.3cm} Prof. Dr. Maximiliam Luppe }\\
			\end{flushleft}
	
			\begin{center}
				\vspace{0.09\textheight}
				\Large{S�o Carlos}\\
				\Large{2018}
			\end{center}
		}
	
\end{titlepage}


%%%%%%%%%%%%%%%%%%%%%%%%%%%%%%%%%%%%%%%%%%%%%%% INSER��O P�GINA EM BRANCO %%%%%%%%%%%%%%%%%%%%%%%%%%%%%%%%%%%%%%%%%%%%%%
\cleardoublepage


%%%%%%%%%%%%%%%%%%%%%%%%%%%%%%%%%%%%%%%%%%%%%%% RESUMO - PORTUGUES %%%%%%%%%%%%%%%%%%%%%%%%%%%%%%%%%%%%%%%%%%%%%%
\
\vspace{0.11\textheight} 

\begin{center}
\textbf{\Huge{Resumo}}
\end{center}

\vspace{0.05\textheight}

Como proposta de projeto para a disciplina SEL0630 - Aplica��o de Microprocessadores II, tem-se o desenvolvimento de um sistema de monitoramento remoto, capaz de automaticamente identificar uma fonte de movimento na imagem e acompanh�-la. Imagens obtidas s�o enviadas para um servidor que pode ser acessado pelo usu�rio do sistema. Ser�o utilizados um {\itshape{BeagleBone black}}, uma {\itshape{webcam}} e um motor controlado por PWM para o projeto, e o {\itshape{software}} ser� desenvolvido em C/C++, com apoio da biblioteca OpenCV.
			
%%%Texto em \underline{\textbf{UM PAR�GRAFO}} apenas - deve conter \underline{tudo} resumidamente (introdu��o, m�todo(s), resultados e conclus�es), de tal forma que seja poss�vel compreender a proposta e o que foi alcan�ado.%%%

\vspace{0.05\textheight}
	
Palavras-Chave: Embarcados, Monitoramento, OpenCV, Remoto.

%%%%%%%%%%%%%%%%%%%%%%%%%%%%%%%%%%%%%%%%%%%%%%% INSER��O P�GINA EM BRANCO %%%%%%%%%%%%%%%%%%%%%%%%%%%%%%%%%%%%%%%%%%%%%%
\cleardoublepage

%%%%%%%%%%%%%%%%%%%%%%%%%%%%%%%%%%%%%%%%%%%%%%% RESUMO - INGL�S %%%%%%%%%%%%%%%%%%%%%%%%%%%%%%%%%%%%%%%%%%%%%%
\
\vspace{0.11\textheight} 

\begin{center}
\textbf{\Huge{Abstract}}
\end{center}

\vspace{0.05\textheight}	
		
As a project proposal for the course SEL0630 - Microprocessors Applications II, we have the development of an intelligent remote monitoring system, capable of automatically identifying, and tracking, a source of movement. Obtained footage is sent to a server, accessible by the system's user. For the project, a {\itshape{BeagleBone black}}, an {\itshape{webcam}} and a PWM controlled motor will be used, and the software will be made in C/C++, with support of OpenCV.

\vspace{0.05\textheight}

Keywords: Embedded, Monitoring, OpenCV, Remote.

%%%%%%%%%%%%%%%%%%%%%%%%%%%%%%%%%%%%%%%%%%%%%%% INSER��O P�GINA EM BRANCO %%%%%%%%%%%%%%%%%%%%%%%%%%%%%%%%%%%%%%%%%%%%%%
\cleardoublepage
%\thispagestyle{empty}
%\newpage
%%%%%%%%%%%%%%%%%%%%%%%%%%%%%%%%%%%%%%%%%%%%%%% RESUMO %%%%%%%%%%%%%%%%%%%%%%%%%%%%%%%%%%%%%%%%%%%%%

%%%%%%%%%%%%%%%%%%%%%%%%%%%%%%%%%%%%% CONFIGURA��ES DOS �NDICES %%%%%%%%%%%%%%%%%%%%%%%%%%%%%%%%%%%%%
%\clearpage
%\thispagestyle{empty}
\listoffigures % �ndice de Figuras
(Se houver...)
\listoftables % �ndice de Tabelas
(Se houver...)
%%%%%%%%%%%%%%%%%%%%%%%%%%%%%%%%%%%%%%%%%%%%%%% INSER��O P�GINA EM BRANCO %%%%%%%%%%%%%%%%%%%%%%%%%%%%%%%%%%%%%%%%%%%%%%
\cleardoublepage

%%%%%%%%%%%%%%%%%%%%%%%%%%%%%%%%%%%%% LISTA DE ABREVIATURAS %%%%%%%%%%%%%%%%%%%%%%%%%%%%%%%%%%%%%
\

\vspace{0.11\textheight} 

\textbf{\Huge{Siglas}}
\vspace{0.05\textheight}
			
\begin{tabbing}
\hspace*{0.5cm}\=\hspace{2.5cm}\= \kill

\> PWM \> \textit{Pulse Width Modulation} - Modula��o de Largura de Pulso \\
% Exemplo de lista de lista de abreviaturas
%\> MVC \> \textit{Model-View-Controller} - Modelo-Vis�o-Controlador \\
%\> POO \> Programa��o Orientada a Objetos \\
%\> UI \>  \textit{User Interface} - Interface do Usu�rio \\
%\> UML \> \textit{Unified Modeling Language} - Linguagem de Modelagem Unificada \\


\end{tabbing}
\cleardoublepage
%%%%%%%%%%%%%%%%%%%%%%%%%%%%%%%%%%%%% CONFIGURA��ES DOS �NDICES %%%%%%%%%%%%%%%%%%%%%%%%%%%%%%%%%%%%%
%\usepackage{fancyhdr}
\pagestyle{fancy}
\fancyhf{} % clear all header and footer fields
\fancyhead[RO, LE] {\thepage}

\fancypagestyle{plain}{\pagestyle{fancy}}

\tableofcontents % �ndice Geral

%%%%%%%%%%%%%%%%%%%%%%%%%%%%%%%%%%%%%%%% ADI��O DOS CAP�TULOS %%%%%%%%%%%%%%%%%%%%%%%%%%%%%%%%%%%%%%%	
\chapter{Introdu��o}
\label{Introducao}

Introdu��o do trabalho.
 
Na reda��o da monografia, que � parte important�ssima do projeto, pois � aquela que ficar� p�blica, precisamos definir muito bem o t�tulo, construir um Resumo com todas as partes de um Resumo, esclarecer o(s) objetivo(s) e construir uma conclus�o completa, tudo isso quase que ao mesmo tempo, pois o trabalho j� terminou.
A Introdu��o deve conter um hist�rico (com muitas refer�ncias atuais-procure nas bases consagradas, como por exemplo IEEE, para apresentar n�o s� refer�ncias vol�teis-aquelas da Internet) do assunto apontando a origem e os avan�os que est�o publicados, ressaltando o "foco" de ataque do projeto.
Depois, compor um bom Embasamento Te�rico citando as refer�ncias de onde est� o assunto todo de cada t�pico (sem se confundir com a apresenta��o dos materiais usados no projeto), pois � o lugar onde estar�o presentes as partes da ci�ncia ou as t�cnicas de forma geral que foram utilizadas para construir a solu��o (exemplo: Sistemas Embarcados � um t�pico e n�o Raspberry Pi, Linux Embarcado � um t�pico e n�o a distribui��o escolhida, por�m Linux � um subt�pico de Sistemas Operacionais). Depois mostrar Material e M�todos, ou Desenvolvimento do Projeto, mostrando logo na entrada do cap�tulo uma figura ou diagrama que apresente de forma geral como as partes est�o relacionadas/conectadas. Mostrar os algoritmos da solu��o em forma de fluxogramas ou usando UML, de maneira clara e completa. Se for necess�rio apresentar trechos de c�digo para ressaltar solu��es ou apresentar abordagens, tamb�m que seja de forma direta e simplificada. C�digos completos dever�o estar nos Ap�ndices para a Monografia final, depois da defesa.  
Assim, consegue-se mostrar o quanto voc� evoluiu com os aprendizados do curso e com aqueles que voc� buscou a mais, e o quanto tem de solu��o legal e atual na sua proposta. 
Valorize as conquistas alcan�adas apresentando e discutindo os resultados com gr�ficos, tabelas, figuras, etc.
Nas Conclus�es, tamb�m n�o esque�a de apontar as defici�ncias ou limita��es que s�o inerentes ao trabalho e que no momento n�o est�o no(s) objetivo(s), e tamb�m as defici�ncias ou limita��es dos procedimentos que vc incorporou ao seu trabalho de outros autores.
Indique as dire��es para os trabalhos futuros, pois o autor conhece como ningu�m as oportunidades de continuidade e avan�os do trabalho.

Exemplo de cita��o de Refer�ncia \cite{referencia1}. Outra refer�ncia para a bibliografia \cite{referencia2}.

Segundo \cite{referencia3} h� uma sequ�ncia l�gica para a reda��o da monografia como apresenta em \cite{enzo}.

Refer�ncia para a figura \ref{logo}.

 \begin{figure}[H]
 	\centering
 	\includegraphics[scale=0.35]{./Resources/proj_sistemas_embarcados.jpg}
 	\caption{Logo da Disciplina.}
 	\label{logo}
 \end{figure}



% % % % % % % % % % % % % % % % % % % % % % % % % % % % % % % % % % % % % % % % % % % % % % % % % % %
\section{Motiva��o}

Descrever a motiva��o do trabalho.


% % % % % % % % % % % % % % % % % % % % % % % % % % % % % % % % % % % % % % % % % % % % % % % % % % %
\section{Objetivo(s)}

Somente o(s) Objetivo(s) do trabalho.


% % % % % % % % % % % % % % % % % % % % % % % % % % % % % % % % % % % % % % % % % % % % % % % % % % %
\section {Justificativas/relev�ncia}

Justificativa do trabalho.


% % % % % % % % % % % % % % % % % % % % % % % % % % % % % % % % % % % % % % % % % % % % % % % % % % %
\section {Organiza��o do Trabalho}

Este trabalho est� distribu�do em XXX cap�tulos, incluindo esta introdu��o, dispostos conforme a descri��o que segue:

Cap�tulo 2: Descreve .....................................................................................

Cap�tulo 3: Discorre sobre .....................................................................................

Cap�tulo 4: Apresenta .....................................................................................




























%\chapter{Especifica��o do Projeto}
\label{Especificacao}

Especifica��o do projeto.


% % % % % % % % % % % % % % % % % % % % % % % % % % % % % % % % % % % % % % % % % % % % % % % % % % %
\section{Se��o 1}

Se��o dentro de um cap�tulo. 



% % % % % % % % % % % % % % % % % % % % % % % % % % % % % % % % % % % % % % % % % % % % % % % % % % %
\section{Se��o 2}

Outra se��o dentro do cap�tulo.


\chapter{Embasamento Te�rico ou Fundamenta��o Te�rica}
\label{EmbasamentoTeorico}

Embasamento te�rico para o desenvolvimento do trabalho.

Leia o texto que est� na Introdu��o e as dicas mais � frente...





\chapter{Material e M�todos ou Desenvolvimento do Projeto}
\label{Material}


% % % % % % % % % % % % % % % % % % % % % % % % % % % % % % % % % % % % % % % % % % % % % % % % % % %
\section{Material}

Material utilizado no projeto.


% % % % % % % % % % % % % % % % % % % % % % % % % % % % % % % % % % % % % % % % % % % % % % % % % % %
\section{M�todos}

M�todos utilizados no projeto.




%%%%%%%%%%%%%%%%%%%%%%%%%%%%%%%%%%%%%% ADI��O DAS BIBLIOGRAFIAS %%%%%%%%%%%%%%%%%%%%%%%%%%%%%%%%%%%%%
\addcontentsline{toc}{chapter}{Refer�ncias}
\renewcommand{\bibname}{Refer�ncias}	
\bibliographystyle{unsrt} % Define o estilo da bibliografia
\bibliography{./Content/References} % Faz referencia ao arquivo ref.bib

%%%%%%%%%%%%%%%%%%%%%%%%%%%%%%%%%%%%%%% T�RMINO DO DOCUMENTO %%%%%%%%%%%%%%%%%%%%%%%%%%%%%%%%%%%%%%%%
\end{document}    
	